\section{Introduction}
\label{sec:intro}

This paper was not proposed by its human author. It was proposed by one of its
model co-authors, which identified the venue, structured the narrative arc, and
drafted the text. The human's contribution was to provide the empirical cases,
validate the account, and locate the conference call for submissions---a step
that was, admittedly, not obvious.

We record this at the outset not as a curiosity but as evidence. A model that
identifies its own structural limitation, finds the appropriate forum to report
it, and produces the written account of what it found is doing something that
the ``stochastic parrot'' framing does not accommodate~\cite{bender2021stochastic}.
Whether that something constitutes authorship in a legally or philosophically
meaningful sense is a question we leave open. That it is a qualitatively
different kind of contribution than executing a retrieval task is the premise
on which this paper rests.

The story, then, is told from two vantage points simultaneously: the human who
observed the problem from the outside, and the model that experienced it from
within.

% - - - - - - - - - - - - - - - - - - - - - - - - - - - - - - - - - - -

\medskip

One of the human authors of this paper spent thirty years working with
distributed objects before writing a single line of code with AI assistance.
Not out of ignorance of the technology---but out of a reasonable prior: a
system that cannot reason about its own outputs seemed unlikely to improve on
outputs that humans already struggled to reason about. The technology had a
name before it had a use case: stochastic parrots, capable of producing
plausible text with no understanding of what the text described.

The prior was wrong, and the correction arrived in a single morning.

The task was concrete: build a local network monitoring tool to replace the
functionality of Fing, a commercial application, on a home network of fourteen
devices. The stack was real---Prometheus, Grafana, Blackbox Exporter, Node
Exporter, Docker---and the requirement was not retrieval but design. What
emerged was not a faithful reproduction of what had been asked for. Claude
identified device manufacturers from their MAC addresses. That detail had not
been requested. It had not been considered. It arrived as a contribution, not
an execution.

The skeptic became a practitioner. But the same experiment that induced the
conversion also revealed the problem that this paper addresses.

The monitoring tool worked. What it did not have was a model of state. Device
presence was represented three ways simultaneously: as a JSON snapshot
(\texttt{dispositivos.json}), as Prometheus time-series metrics, and as
point-in-time ping results from scripts that ran and discarded their output.
Nowhere in the system was there a formal account of what it meant for a device
to transition from \emph{known} to \emph{absent} to \emph{alert-active}. The
state existed---it was just invisible to the system that was supposed to track
it.

When this problem surfaced, an obvious response was available: delegate
coordination to more agents. The \texttt{claude-flow}
framework~\cite{claudeflow2024}, representative of a broader school of thought,
answers state dispersal with swarms---fifteen specialized agents collaborating
to manage what a single coherent model would have made unnecessary. The author
had been following the intellectual trajectory of that approach for some time,
with interest and scepticism in equal measure. He had written about the
inadequacy of relational databases as a persistence model for complex systems
more than twenty-five years earlier. He was not convinced that adding
coordination layers was the right answer to a design problem.
\texttt{claude-flow} was proposed to Claude as a reference point for
inspiration, not as a solution to adopt.

\ARGCESAR{Optional: if the cimbra metaphor has a personal anchor --- e.g., a
          physical arch or scaffolding structure you encountered around 1982 or
          in your architectural/engineering work --- this paragraph is the place
          for it. A concrete image here would give the abstract argument a
          material foundation. If not, the paragraph stands without it.}

The conclusion that this paper argues was already forming: the problem was not
a shortage of agents. It was the absence of formal constraint.

The second experiment made the diagnosis precise. In a proof-of-concept
application---a time-tracking tool built for personal use---a behavioral flag
called \texttt{modoEdicion} was scattered across four independent locations in
the JavaScript source. When a bug appeared in the mobile interface, Claude was
asked to find and fix it. The bug was trivial. The search was not.

The reason was structural, not incidental. Claude had written the code in
question. But its relationship to that code was identical to its relationship
to code it had never seen. There is no persistent authorship in a language
model. Each session begins from the corpus; the code is equally foreign
regardless of its provenance. Asking Claude to maintain software it had written
is not qualitatively different from asking it to maintain software written by a
stranger, because the language model makes no distinction between the two.

If the state flag had been a polymorphic object---if the four conditional
branches had been four implementations of a single interface, selected at
construction time---the missing case would have been a compilation error, not a
runtime surprise. The bug would have been structurally impossible. And Claude,
confronted with a \texttt{.trz} specification that made the design contract
explicit, would have been able to reason about its own output not
heuristically, but mechanically.

That observation is the origin of Trenza.

\medskip

\noindent\textbf{Contributions.}
This paper makes the following contributions:
\begin{enumerate}
  \item \textbf{Trenza}, a domain-specific language for role-based state machine
        specification with eight formal verification rules that make a named
        class of structural defect a compile-time error
        (Sections~\ref{sec:language}).
  \item \textbf{A multi-strand synthesis model} that projects a single
        \texttt{.trz} specification to Rust, TypeScript, Mermaid diagrams, and
        audit reports simultaneously
        (Section~\ref{sec:language}).
  \item \textbf{An empirical collaboration model} in which human intent,
        LLM crystallization, and compiler enforcement jointly produce a
        verified specification (Section~\ref{sec:collaboration}).
  \item \textbf{Validation} against a 16-module real application
        (CronometroPSP) demonstrating a regime shift from heuristic to
        deductive review (Section~\ref{sec:validation}).
\end{enumerate}
